
\chapter{Эмитент и~эквайер: две стороны одной транзакции}\label{ch:issuer-acquirer}

Один короткий жест покупателя у~терминала или на странице оплаты
запускает два разных процесса. Эмитент за секунды проверяет доверие:
действительна ли карта, есть ли деньги, не слишком ли велик риск
мошенничества. Эквайер решает другую задачу: принять запрос, провести
его по нужному маршруту и корректно рассчитаться с~продавцом.

\section{Эмитент}\label{sec:ia-issuer}

Эмитент (issuer) выпускает карту и принимает на себя кредитный
риск транзакции. В~момент покупки он должен решить, готов ли
оплатить операцию через карточную сеть, а~затем взыскать средства
с~держателя карты: списать с~дебетового счёта или увеличить
задолженность по кредитному. Ответ за секунды, а~не за дни,
держится на наборе систем, каждая из которых отвечает за свой
слой реальности: состояние карты, правила авторизации,
аутентификацию, токены и претензии после операции.

\subsection*{Система управления картой (CMS)}

Система управления картой (Card Management System, CMS) ведёт
жизненный цикл карты: выпуск, смена статуса, блокировка, перевыпуск.
CMS хранит реквизиты и связь со счётом, поддерживает статусы и~лимиты,
формирует задания на персонализацию и фиксирует события, которые
позже потребуются поддержке, антифроду и разбору диспутов.

На этапе персонализации на карту загружаются EMV-апплет, ключи
эмитента (IMK), параметры управления риском (порог офлайн-авторизации
floor limit и~офлайн-лимиты), а~также данные платёжного приложения
(AID, Application Label) (см. главу~\ref{ch:emv})~\cite{emvco-book3,emvco-book4}.
Для карт Мир персонализация дополнительно включает установку
приложения Мир (МПА) с~поддержкой российской криптографии (ГОСТ);
конкретный криптографический профиль зависит от программы эмиссии
и продолжающегося перехода НСПК на отечественные алгоритмы
(см. главу~\ref{ch:mir}).

\subsection*{Авторизационный хост}

Авторизационный хост решает, можно ли по карте провести конкретную
операцию. Когда запрос доходит от карточной сети до эмитента, хост
последовательно снимает несколько рисков: проверяет статус карты и
срок её действия, смотрит доступность средств или свободного
кредитного лимита, сверяет лимитный контур, оценивает признаки
нетипичного поведения (см. главу~\ref{ch:fraud}) и~для EMV-транзакций
верифицирует криптограмму ARQC через HSM~\cite{emvco-book2,emvco-book3}.

Если все проверки пройдены, хост формирует ответ \texttt{00}
(Approved) и генерирует ARPC (Authorization Response Cryptogram),
который возвращается карте для взаимной аутентификации.

Граничный случай: CMS недоступен, авторизационный хост работает.
Без CMS хост не видит ни актуального статуса карты (заблокирована?
перевыпущена?), ни текущего лимита, и~у~него остаётся два варианта.
Первый -- отклонить все транзакции: безопасно, но парализует весь
карточный портфель. Второй -- работать по кэшу: хост хранит копию
критических полей (статус, лимит, порог офлайн-авторизации) с TTL
5--15~мин и продолжает авторизовать в~пределах кэшированных лимитов.
Транзакции сверх кэшированного лимита отклоняются; для остальных
эмитент принимает риск одобрения по устаревшим данным. После
восстановления CMS хост синхронизирует остатки и фиксирует
расхождения. На портфеле в~миллионы карт разница между 5~мин и
30~мин простоя CMS оборачивается сотнями одобренных транзакций
по заблокированным картам~\cite{emvco-book3}. Если хотя бы одна
проверка не проходит, хост возвращает соответствующий response code
(\texttt{51}: Insufficient Funds, \texttt{05}: Do Not Honor,
\texttt{14}: Invalid Card Number и~так далее), и транзакция
отклоняется. Решение собирается из цепочки проверок: каждая
отсекает часть риска, пока не остаётся либо одобрение, либо
объяснимый отказ~\cite{emvco-book3}.

\subsection*{Сценарии эмитента (issuer scripting)}

Решение эмитента не всегда сводится к Approved или Declined --
иногда вместе с~ответом нужно изменить поведение самой карты.
Для этого EMV-стандарт предусматривает issuer scripting --
механизм передачи команд карте через авторизационный ответ.
Script 71 выполняется до генерации финальной криптограммы
(TC или AAC), Script 72 -- после.~\cite{emvco-book3} На практике
issuer scripting используют, когда нужно обновить офлайн PIN без
визита в~банк, сбросить офлайн-счётчики после успешной онлайн-авторизации,
заблокировать приложение прямо во время контакта карты с~терминалом
или обновить параметры управления риском, например лимиты на
офлайн-транзакции.

Issuer scripting сшивает авторизацию и жизненный цикл карты в~одну
операцию, но применять его получается не везде: терминалы по-разному
надёжно обрабатывают скрипты, а~у~бесконтактной транзакции окно
слишком короткое, и~сложное обновление просто не успевает
завершиться (см. главу~\ref{ch:contactless}).

\subsection*{ACS: сервер аутентификации}

В~интернет-платежах одной проверки карты, лимита и криптограммы
эмитенту недостаточно: нужно ещё понять, действительно ли покупку
инициировал держатель. Эту задачу решает ACS -- компонент эмитента,
отвечающий за аутентификацию 3-D Secure (см. главу~\ref{ch:3ds}).
Когда продавец инициирует проверку 3-D Secure, запрос приходит в
ACS, и~тот решает, пропустить ли операцию по сценарию frictionless,
без участия держателя, или потребовать challenge -- одноразовый
код, биометрию или подтверждение в~мобильном
приложении.~\cite{emvco-3ds,emvco-3ds-whitepaper}

Решение ACS строится на Risk-Based Authentication (RBA): учитываются
устройство, IP-адрес, история прошлых транзакций, сумма, профиль ТСП
и~другие признаки контекста. Качество этого решения напрямую
определяет конверсию. Слишком строгий ACS отправляет на challenge
операции, которые можно пропустить без трения, и~продавец теряет
покупателей. Слишком мягкий ACS пропускает мошенничество и
увеличивает убытки эмитента.~\cite{emvco-3ds-whitepaper}

\subsection*{Управление токенами}

После выпуска карты её реквизиты живут на пластике, в~телефоне,
в~браузере и~у~продавца в~виде сетевых токенов (DPAN)
(см. главу~\ref{ch:tokenization}). Эмитент участвует в~этом
жизненном цикле с~самого начала: получает от провайдера токенизации
(TSP) запрос на токенизацию, принимает решение об одобрении после
процедуры ID\&V (identification and verification) держателя, а
затем управляет состоянием токена -- активирует, приостанавливает
или удаляет.~\cite{emvco-payment-tokenisation,visa-vts,nspk-tsp}

Токен живёт как отражение исходной карты. Если держатель блокирует
карту, эмитент обязан приостановить и связанные с~ней токены. Если
карта перевыпущена, нужно обновить привязку токенов к~новому PAN.
Иначе платёжный опыт распадается: держатель видит в~приложении
банка новую карту, а~у~кошелька или продавца ещё живёт старое
представление о~ней.~\cite{emvco-payment-tokenisation,visa-vts}

\subsection*{Работа с~диспутами}

Держатель карты с~претензией обращается не в~карточную сеть и~не
к~эквайеру, а~к~эмитенту -- своему банку: \enquote{я не совершал
эту покупку} или \enquote{товар не получен}. Эмитент проверяет
обоснованность жалобы и~при наличии оснований инициирует чарджбэк
через карточную сеть, направляя его эквайеру
(см. главу~\ref{ch:disputes}).~\cite{visa-core-rules,mc-chargeback-guide,visa-vcr-merchants}

На практике контур начинается с~приёма и классификации обращения.
Дальше эмитент сопоставляет ситуацию с~правилами сети: есть ли
подходящий код причины, не пропущены ли сроки, достаточно ли
документов и соблюдены ли обязательные условия для чарджбэка.
После этого эмитент формирует чарджбэк-сообщение, а~если эквайер
присылает повторное представление с compelling evidence, эмитент
решает, принять объяснение или идти дальше, к pre-arbitration.
Разбор диспутов -- продолжение авторизации задним числом:
авторизация отвечает на вопрос \enquote{можно ли пустить деньги},
диспут -- на вопрос \enquote{надо ли это решение
пересмотреть}.~\cite{visa-core-rules,mc-chargeback-guide,visa-vcr-merchants}

\section{Эквайер}\label{sec:ia-acquirer}

Эквайер (acquirer) видит ту же операцию с~противоположной стороны.
Он подключает терминал или платёжный шлюз, доставляет авторизационный
запрос в~карточную сеть, принимает расчёт от эмитентов и перечисляет
деньги продавцу за вычетом комиссий. Эквайерская архитектура покрывает
приём платежа, удержания, выплату продавцу и~правила расчёта.

\subsection*{Система управления терминалами (TMS)}

\definitionnote{TMS}{Система удалённого управления парком
  POS-терминалов: обновления, ключи, конфигурация терминала,
мониторинг состояния и диагностика неисправностей.}

В~офлайн-эквайринге первый физический слой -- терминал. У~крупного
эквайера TMS управляет десятками и сотнями тысяч устройств, и~главная
задача здесь -- постоянно держать актуальную конфигурацию: обновлённые
BIN-таблицы, свежие ключи, корректные EMV-параметры.

Без удалённой инъекции ключей (Remote Key Injection, RKI) эквайринг
плохо масштабируется: ручная церемония загрузки ключей на каждый
терминал быстро упирается в~операционный потолок. RKI снимает это
ограничение -- терминалы разворачиваются и обновляются без физического
доступа.~\cite{visa-tpa-registration-program-faqs}

\subsection*{Управляющий контур}

Для интернет-эквайринга и омниканальных сценариев одних терминалов
мало. Нужен ещё слой, в~котором живут профили ТСП, конфигурация
MID/TID, правила дескрипторов, маршрутизация по процессорам,
параметры 3-D Secure, лимиты на возвраты и~выплаты, а~также связка
с~платёжным шлюзом из главы~\ref{ch:gateway}. Этот слой и~называют
управляющим контуром.

Управляющий контур определяет, какой MCC и~какой дескриптор увидит
эмитент, через какой BIN эквайера пойдёт транзакция и в~какой контур
мониторинга попадёт ТСП. Если такие решения расползаются между
таблицами, тикетами и~ручными настройками в~кабинете процессора,
эквайер теряет управляемость: одна и~та же команда ТСП начинает
по-разному выглядеть для сети, для команды риска и~для финансовой
сверки.~\cite{visa-merchant-data-standards-manual}

У~крупных эквайеров в~управляющем контуре сходятся подключение ТСП,
маршрутизация между несколькими процессорами и операционная политика.
Здесь решается, кому дать прямой маршрут на конкретного процессора,
кому оставить один канал, а~кому включить усиленный holdback или
ручную проверку выплат.

\subsection*{Подключение ТСП}

Эквайер не может включить приём платежей по одной заявке с~сайта.
Прежде чем дать ТСП доступ к~карточной сети, нужно понять, кто
перед ним, какой бизнес он ведёт и~какой риск несёт. Эту задачу
решает контур подключения ТСП -- набор проверок и~настроек,
задающих и~допуск, и экономику отношений.

\textbf{KYB (Know Your Business)}: проверка юридического лица --
регистрация, бенефициары, финансовое состояние, репутация. Эквайер
обязан убедиться, что ТСП не занимается запрещённой деятельностью,
не находится в~санкционных списках и~ведёт легитимный бизнес.
Глубина проверки зависит от уровня риска: интернет-казино проверяют
строже, чем продуктовый магазин.

\textbf{MCC-классификация}: присвоение четырёхзначного кода
категории (Merchant Category Code, MCC), задающего тип бизнеса.
От MCC зависит ставка interchange, правила 3-D Secure, лимиты
чарджбэков и иногда сама допустимость операции. MCC 7995, например, охватывает ставки, лотереи и~казино,
поэтому для таких операций сети и~эквайеры вводят отдельные
ограничения и усиленный
риск-контроль.~\cite{visa-merchant-data-standards-manual,visa-payment-facilitator-marketplace-risk-guide}
В~России онлайн-казино запрещены, лотерея -- госмонополия, а
расчёты по ставкам легальных букмекерских контор и тотализаторов
идут не через обычный карточный эквайринг, а~через единый центр
учёта переводов интерактивных ставок (ЕЦУПС, оператор -- НКО
\enquote{Мобильная Карта}) на основании статьи~14.2 244-ФЗ и
Распоряжения Президента РФ от 25~августа 2021~года № 236-рп
о~наделении НКО \enquote{Мобильная Карта} статусом единого центра
учёта~\cite{fz-244,etsups}.

\practicenote{%
  Неправильный MCC -- одна из самых частых ошибок при подключении,
  и последствия у~неё вполне материальные. ТСП платит неверный
  interchange, транзакции некорректно проходят антифрод-фильтры,
  а~при аудите карточной сети ошибка оборачивается штрафом и
  требованием рекатегоризации. Выбор MCC требует той же
  внимательности, что и~проверка юридических документов.
}

\textbf{Underwriting}: оценка финансового риска ТСП. Эквайер
смотрит, какой объём транзакций ожидается, каков средний чек,
насколько вероятны чарджбэки и какой кредитный лимит можно
предоставить до расчёта. Underwriting задаёт условия, на которых
ТСП живёт в~системе: ставку комиссии, размер holdback, наличие
rolling reserve и~график
выплат.~\cite{visa-payment-facilitator-marketplace-risk-guide}
\subsection*{Мониторинг ТСП}

После подключения риск никуда не девается, и~эквайер обязан
непрерывно наблюдать за ТСП в~рабочем режиме. Мониторинг опирается
на несколько групп сигналов:

\begin{itemize}
  \item \textbf{Доля чарджбэков (chargeback ratio)}: отношение числа чарджбэков
    к~числу транзакций. Карточные сети устанавливают пороги
    и программы комплаенса (для Visa сейчас это единый
      VAMP, Visa Acquirer Monitoring Program -- метрика на уровне
    эквайерского портфеля, объединившая мониторинг фрода и~споров;
    у Mastercard -- Excessive Chargeback Program, ECP,
    и Excessive Fraud Merchant, EFM)
    (см. главу~\ref{ch:disputes});
    их превышение ведёт к~штрафам, усиленному мониторингу
    и в~крайнем случае к~отключению
    ТСП.~\cite{visa-vamp-intro,mc-compliance-programs}
  \item \textbf{Fraud rate}: доля мошеннических транзакций.
  \item \textbf{Доля возвратов (refund ratio)}: аномально высокий процент
    возвратов, который может указывать либо на проблемы
    с~товаром, либо на мошенническую схему
    (triangulation fraud).
  \item \textbf{Velocity anomalies}: резкий рост объёма
    транзакций, изменение среднего чека или появление
    нехарактерных BIN-диапазонов.
\end{itemize}

Помимо собственного мониторинга, российские операторы карточного
контура обязаны интегрироваться с~автоматизированной системой
обработки инцидентов ФинЦЕРТ Банка России (АСОИ ФинЦЕРТ): Указание
Банка России от 8~октября 2018~года № 4926-У~закрепляет порядок и
формы информационного обмена по инцидентам в~платёжных системах.
Через АСОИ участники получают предупреждения о~компрометации
реквизитов и сами сообщают о~выявленных операциях без согласия
клиента.

Добросовестное ТСП редко становится проблемным за один день, но
почти всегда заранее оставляет следы в~метриках: внезапный рост
объёма, смещение географии, неестественные возвраты, ухудшение
доли одобрений.~\cite{visa-payment-facilitator-marketplace-risk-guide}

\importantnote{%
  Ответственность за ТСП лежит на эквайере. Если ТСП оказался
  мошенником или нарушил правила карточной сети, штрафы и
  финансовые потери несёт именно эквайер, а~не сеть и~не эмитент.
  Мониторинг ТСП защищает от прямых финансовых потерь: пропущенный
  сценарий bust-out, когда ТСП сначала наращивает объём, а~затем
  исчезает с~деньгами, обходится эквайеру в миллионы.~\cite{visa-payment-facilitator-marketplace-risk-guide,visa-acceptance-entities}

}
\subsection*{Ответственность за расчёт}

Для продавца эквайринг заканчивается в~момент поступления денег.
Эквайер собирает клиринг, удерживает комиссии, применяет rolling
reserve и перечисляет средства продавцу.

Эквайер обязан прозрачно объяснить каждую выплату: какая сумма
удержана как комиссия, какая отправлена в~резерв, какая зачтена
против возвратов и в~какой валюте считается обязательство. Кроме
того, на нём риск late presentment, временных разрывов по чарджбэкам
и кассовых разрывов между получением денег от сети и~выплатой
продавцу. Эквайер же задаёт политику выплат: T+1, T+2, еженедельные
выплаты, отложенный расчёт для high-risk MCC -- правила, которые
для продавца нередко важнее самой доли одобрений.

Holdback и rolling reserve покрывают временной разрыв между
авторизацией и закрытием окна возможных чарджбэков, а~до закрытия
этого окна может пройти несколько месяцев. Если продавец за это
время прекратил деятельность или не может покрыть возвраты, потери
падают на эквайера. Rolling reserve -- буфер, из которого эквайер
покрывает эти потери. Размер резерва отражает оценку риска ТСП:
высокая доля чарджбэков, сезонный бизнес или новое ТСП без истории
получают более высокий процент удержания.

\section{Платёжный фасилитатор (PayFac) и~модель подмерчантов}\label{sec:ia-payfac}

Классическая модель эквайринга работает, пока ТСП немного и~каждое
подключается как отдельный клиент банка. Если же платформа хочет
быстро вывести в~платёжный контур сотни или тысячи мелких продавцов
(маркетплейс или SaaS-платформа с~большим числом контрагентов),
банковская процедура становится слишком медленной. Индивидуальное
подключение каждого продавца через банк занимает дни или недели и
начинает тормозить саму бизнес-модель.

\definitionnote{PayFac}{Агент спонсирующего эквайера, который
  подключает подмерчантов от его имени, получает расчёт в~их пользу и
отвечает за соблюдение правил сети и регуляторных требований.}

В~модели PayFac платформа регистрируется спонсирующим эквайером
в~карточной сети как Payment Facilitator и получает право подключать
sponsored merchants по ускоренной схеме. У~подмерчанта может быть
собственный идентификатор, но договорная и расчётная рамка идёт
через связку PayFac + его эквайер: эквайер перечисляет средства
PayFac, а~тот уже рассчитывается с~подмерчантами. Главный выигрыш --
скорость: нового продавца подключают за минуты, а~не за недели.
PayFac автоматизирует KYB, присваивает MCC и~начинает приём платежей,
не дожидаясь полного банковского
цикла.~\cite{visa-acceptance-entities,visa-payment-facilitator-marketplace-risk-guide}

Скорость покупается концентрацией ответственности. PayFac отвечает
перед карточной сетью и эквайером за своих подмерчантов. Если один
из них оказался мошенником, штрафы ложатся на PayFac. Если у
подмерчанта выросла доля чарджбэков, отвечает тоже PayFac. Поэтому
PayFac строит собственные процессы KYB, мониторинга ТСП и~работы
с~диспутами, по строгости мало отличающиеся от
банковских.~\cite{visa-acceptance-entities,visa-payment-facilitator-marketplace-risk-guide}

Главный структурный риск PayFac-модели в~том, что подмерчанты
находятся вне прямой видимости карточной схемы. Схема видит PayFac
как единый субъект; что происходит внутри его портфеля, ей напрямую
не известно. Если подмерчант нарушает правила, схема замечает это
только через отклонения в~агрегированных показателях PayFac.
Ответственность за обнаружение и финансовые последствия одинаково
лежат на PayFac.

\importantnote{%
  PayFac и ISO (Independent Sales Organization) -- разные роли.
  ISO привлекает и обслуживает ТСП для банка-эквайера, получая
  агентское вознаграждение без прямого участия в~расчёте. PayFac
  стоит прямо в~денежной цепочке: эквайер перечисляет средства ему,
  а~он затем рассчитывается с~подмерчантами. Эта разница задаёт и
  масштаб ответственности, и~набор регуляторных
  требований.~\cite{visa-tpa-registration-program-faqs,visa-acceptance-entities}

}
\subsection*{Когда выбирать PayFac}

PayFac оправдан, когда платформе нужно быстро подключать большое
число небольших ТСП. Маркетплейсы и подписочные платформы выигрывают
от того, что сами управляют сценарием подключения, удерживают свою
комиссию и контролируют выплаты. Там, где ТСП мало и~они крупные,
прямой эквайринг чаще оказывается проще и~дешевле.

\textbf{PayFac} (Payment Facilitator): платформа регистрируется в
сети через спонсирующего эквайера, ведёт KYB подмерчантов, несёт
ответственность за чарджбэки и контролирует выплаты. Подходит при
>100 подмерчантов и потребности в~быстром подключении.
\textbf{MoR} (Merchant of Record): платформа -- единственный ТСП в
глазах сети, а подмерчанты -- её внутренняя абстракция. Проще
регуляторно, но вся ответственность (чарджбэки, налоговые
обязательства, возвраты) замыкается на платформу. Подходит для SaaS
с~единым продуктом. \textbf{Pass-through} (прямой эквайринг): каждое
ТСП заключает договор с~эквайером напрямую, а платформа только
маршрутизирует трафик. Ответственность платформы минимальна, но
подключение каждого ТСП занимает дольше, а управлять выплатами
платформа не
может~\cite{visa-acceptance-entities}. На практике выбор между
pass-through, MoR и PayFac определяется не числом ТСП, а сочетанием
факторов: однородность продукта, готовность платформы к KYB и
контролю выплат, готовность взять на себя риски подмерчантов и
распределение ответственности по чарджбэкам.

Цена этого удобства:
\begin{itemize}
  \item повышенная ответственность за мошенничество и чарджбэки
    подмерчантов;
  \item необходимость регистрации через спонсирующего эквайера
    в~карточных сетях;
  \item инвестиции в~собственные процессы KYB и AML;
  \item в~некоторых юрисдикциях возникает необходимость дополнительной
    лицензии или специального статуса.
\end{itemize}

PayFac несёт эти обязанности, потому что встроен в~модель спонсорства
эквайера и действует по правилам карточной
сети.~\cite{visa-acceptance-entities,visa-payment-facilitator-marketplace-risk-guide}

\practicenote{%
  Если вы строите платформу и выбираете между PayFac и~прямым
  эквайрингом для каждого ТСП, оцените, сколько ТСП реально
  планируете подключить в~первый год. На десятках прямой эквайринг
  обычно проще и~дешевле. На сотнях и~тысячах PayFac окупается за
  счёт скорости подключения, но требует серьёзных инвестиций в~риск,
  соблюдение требований и операционные процессы. Промежуточный
  вариант -- готовое решение для маркетплейсов и платформ у~крупного
  российского эквайера или PSP (например, ЮKassa, Т-Бизнес, Сбер
  для платформ): скорость подключения покупается вместе с~чужой
  инфраструктурой и готовыми процедурами KYB.

}

\subsection*{Российская параллель: банковский платёжный агент}

PayFac-модель родилась внутри правил карточных сетей и применяется
там же. Российский правопорядок оперирует другой рамкой --
институтом банковского платёжного агента (БПА) и банковского
платёжного субагента, закреплённым статьями~14 и 14.2
161-ФЗ~\cite{fz-161}; для БПА, осуществляющих операции платёжного
агрегатора, дополнительно применяется Указание Банка России от
06.04.2020 № 5429-У~\cite{cbr-5429u}. Кредитная организация вправе
привлекать БПА для приёма наличных, переводов, идентификации и
иных операций в~случаях и~в~пределах, установленных законом.

Для платформы, подключающей подмерчантов и работающей с~российским
эквайрингом, БПА -- не прямая замена PayFac из правил Visa/Mastercard,
а отдельный регуляторный режим с~собственным набором обязанностей:
для БПА -- операторов платёжного агрегатора предусмотрено включение
в~перечень, который ведёт Банк России (общего публичного реестра
всех БПА не существует); выполнение требований 115-ФЗ по
идентификации и контролю операций (глава~\ref{ch:compliance-aml});
ведение специального банковского счёта, на который зачисляются
принимаемые средства и~с~которого они перечисляются; ограничения
по перечню операций, которые БПА имеет право выполнять от имени
банка. Банковский платёжный субагент находится на следующем уровне
договорных отношений и привлекается уже самим БПА с~дополнительными
ограничениями.

На практике у~российской платформы выбор идёт не между \enquote{Stripe
Connect} и \enquote{Adyen for Platforms}, а~между несколькими
вариантами на российском рынке: договор БПА с~одним или несколькими
банками; интеграция через PSP, который сам уже имеет статус БПА и
предлагает решение для маркетплейсов как собственный продукт;
работа в~карточном контуре в~роли спонсируемого PayFac у~российского
эквайера -- участника системы \enquote{Мир} (если такая модель
доступна по правилам НСПК для конкретной категории платформ).
Архитектуру выплат и~сверки нужно проектировать с~учётом того, какие
именно режимы задействованы; ссылка на универсальные PayFac-правила
без оговорки про БПА для российского рынка неточна.
